\subsection[MORE THAN ONE UNIT PER HEX]{\mbox{MORE THAN ONE UNIT PER HEX}}

\begin{enumerate}
  \item Up to thirty Stacking Factors may be present in any one hex at any one time. This limit can never be exceeded, even if units must be eliminated from the game to avoid exceeding it.

  \item Each type of unit counter has a specific Stacking Factor (as listed on the MOVEMENT \& TERRAIN EFFECTS CHART). Any combination of units may be placed together in a hex, as long as their combined Stacking Factor does not exceed thirty.

  \item A unit counter or Stack of units cannot be moved into a hex containing other friendly units without Stacking them together. A unit counter or Stack of units cannot be moved out of a hex containing other friendly units without Unstacking them. Movement Allowance Factors are always expended when any Stacking or Unstacking takes place in a hex.

  \item In no case may units be Stacked with enemy units in the same hex.

  \item Unit counters with various Combat Factors are provided for the player's convenience, to avoid large Stacks when strong forces are present in one hex. Large Combat Factor units can replace a number of smaller Combat Factor units of the same type, or vice-versa, at any time a player wishes to do so, as long as the number of Combat Factors remains the same. In other words, treat the counters as if they were "change" - i.e. five pennies (five "1" Combat Factor unit counters) equals one nickel (one "5" Combat Factor unit counter).
\end{enumerate}

\textbf{Example of Play \#2: Stacking Factors}

The Stack of thrity Infantry Combat Factors does not exceed the limits for their hex, as Infantry Combat Factors count for only one Stacking Factor each. The Stack of fifteen Artillery counters is illegal, because Artillery Combat Factors count for three Stacking Factors each. Thus, the Artillery Stack exceeds the limit by fifteen Stacking Factors - no more than ten Artillery Combat Factors could ever be in one hex.

TODO: Insert image here.

\subsubsection[STACKING \& UNSTACKING UNITS]{\mbox{STACKING \& UNSTACKING UNITS}}

\begin{enumerate}
  \item The placing of units in a stack with other units in the same hex, or taking a unit or units from a hex where it was previously stacked with other units represents a reorganization of the units in the hex, some marching and counter-marching, etc, all of which takes time and costs Movement Allowance Factors.

  \item Stacking or Unstacking a unit or units costs a varying number of Movement Allowance Factors for \textit{all} units in the hex, depending on the terrain in the hex in which it is done. These costs are shown on the MOVEMENT \& TERRAIN EFFECTS CHART.

  \item In hexes where there are several types of terrain features, the highest cost is the one used.

  \textit{For example, a unit in a hex that contains both clear and slope terrain will use up the number of Movement Allowance Factors appropriate for slope terrain when Stacking and Unstacking units.}

  \item Roads are ignored as terrain for the process of Stacking and Unstacking units.

  \item Stacking of units may be done at any time during a Player's turn, provided that each involved unit has enough Movement Allowance Factors left to do so. Similarly, units may be Unstacked at any time during a players Turn, provided that each involved unit has enough remaining Movement Allowance Factors to do so.

  \item Once stacked together, the units in the hex may continue to move together as if they were all one unit. They cannot split up until Movement Allowance Factors are used to Unstack them.

  \item All units in the same hex must be Stack together - there cannot be several different Stacks in the same hex.
\end{enumerate}

\textbf{Example of Play \#3: Stacking \& Unstacking}

TODO: Insert image here.

Note that the Cavalry unit expends 2 MAF to Unstack from the 9-6 unit since this is a slope hex. The Cavalry unit must expend 1 MAF to stack with the 6-6, then expend 1 MAF to Unstack with it, in order to move through the hex that the 6-6 occupies. A total of 9 MAF are expended by the Cavalry unit for this move. Note that Infantry units in this example were also required to expend MAF for Stacking and Unstacking.

\subsubsection{STACKING COUNTERS AND STACKING BOXES}

\begin{enumerate}
  \item A large Stack of many unit counters in the same hex is awkward to move, and difficult for a player to handle. For this reason, the Stacking Counters and Stacking Boxes are provided (the Stacking Boxes also have a purpose for Combat, explained later in the rules).

  \item Each Stacking Counter is identified by a letter, and for each Stacking Counter there is also a Stacking Box printed on the Mapboard that has a corresponding letter (there are also Stacking Boxes without an identifying letter, whose use is explained later in the rules).

  \item For convenience, players may at any time remove a Stack of unit counters from a hex on the Mapboard. A Stacking Counter is then placed in that hex, and the unit counters are placed in the Stacking Box with the letter corresponding to the letter on the Stacking Counter. The Stacking Counter now represents that Stack of units on the Mapboard hex, and the player can spread out his units in the Stacking Box for greater ease of play and/or Combat purposes.
\end{enumerate}